% Generated from tex/chapters/ch04/11_発展的な読み物.tex for the SWEBOK structure tree
\subsection{発展的な読み物}

\par\smallskip\noindent\refstepcounter{table}\label{tab:ch04-05}\textbf{表 \thetable: 第04章 ソフトウェア構築 表05}\par\smallskip
\begin{longtable}{L{0.34\linewidth}L{0.58\linewidth}}
\toprule
\textbf{文献} & \textbf{読むと得られること} \\
\midrule
McConnell, \textit{Code Complete} & コーディング、複雑さの管理、構築品質、レビュー、テスト、チューニングを実務的に学べる。構築を初めて体系的に学ぶときの中心文献である。 \\
Sommerville, \textit{Software Engineering} & ソフトウェア工学全体の中で、構築、再利用、分散システム、テストを位置づけて理解できる。 \\
Kim et al., \textit{The DevOps Handbook} & 継続的デリバリ、DevOps、フィードバックループ、運用からの学習を学べる。 \\
Fowler and Beck, \textit{Refactoring} & 既存コードの設計を改善し、変更しやすくする具体的な方法を学べる。 \\
Robert C. Martin, \textit{Clean Code} & 読みやすく保守しやすいコードを書くための原則を学べる。 \\
IEEE Std. 1517-2010 & 再利用に基づく開発、再利用のための構築、再利用による構築のプロセスを学べる。 \\
ISO/IEC 12207 & ソフトウェアライフサイクルプロセス全体の中で、構築、統合、再利用を位置づけられる。 \\
\bottomrule
\end{longtable}
